\begin{figure}[htbp]
  \centering
  \begin{tikzpicture}[x=1.1cm,y=0.7cm]
    \node[anchor=east,font=\small] at (0,1) {Cycle};
    \foreach \cycle in {1,...,8} {
      \node[font=\small] at (\cycle,1) {\cycle};
    }
    \foreach \instruction in {1,...,4} {
      \node[anchor=east,font=\small] at (0,1-\instruction) {$I_{\instruction}$};
      \foreach \stage/\stageindex in {IF/0,ID/1,EX/2,MEM/3,WB/4} {
        \node[draw=bookaccent,fill=booklight,minimum width=1.0cm,minimum height=0.6cm,font=\scriptsize]
          at (\instruction+\stageindex,1-\instruction) {\stage};
      }
    }
    \draw[bookwarm,thick,rounded corners] (4.5,0.48) rectangle (5.5,-3.48);
    \node[font=\small,align=center] at (4.5,-4.5) {After filling, one instruction finishes each cycle.\\Each instruction still crosses five stages.};
  \end{tikzpicture}
  \caption{Four instructions take eight cycles in an ideal five-stage pipeline. Highlighted cycle 5 contains work from all four instructions. Assume one cycle per stage, no hazards, and separate resources wherever stages overlap.}
  \label{fig:five-stage-overlap}
\end{figure}
